\section{Introduction}
\label{sec:introduction}

High-resolution orbital imagery from the Lunar Reconnaissance Orbiter (LRO) now covers the Moon at resolutions down to half a metre per pixel, far beyond manual analysis capacity \citep{lroc_usgs}. Impact craters, wrinkle ridges, lobate scarps, and lava-tube skylights encode the thermal, tectonic, and impact history of the lunar surface, and their automated identification is a prerequisite for geological mapping and the selection of future landing sites \citep{pia12954,nasa2023southpole}. The computer-vision literature offers several distinct formulations of this task --- semantic segmentation (a label per pixel), instance segmentation (individual objects with masks), single-stage object detection (bounding boxes), and classical image processing --- and it is not obvious a priori which suits planetary terrain: published comparisons are typically confounded by differences in datasets, train/test splits, metrics, and operating thresholds.

This work removes those confounds by comparing one representative of each family --- a thresholding-and-morphology baseline, a compact U-Net (SmallUNet), an adapted Mask R-CNN, and YOLOv8 --- on the \emph{same} data, the \emph{same} leakage-free spatial split, and the \emph{same} metrics. The comparison is organised per feature class rather than per architecture, because the label channels range from an impact-crater layer covering 82\% of all pixels to classes with fewer than one positive pixel in $10^{4}$, and no single aggregate ranking is meaningful across that range. Two methodological experiments frame it. First, we show that the saturated crater channel reduces naive scores to a base-rate artefact --- a predict-everything baseline attains F1 $=0.90$ --- and adopt prevalence-robust metrics accordingly. Second, we quantify the train--validation leakage induced by 50\%-overlapping tiles, by retraining identical detectors under a random tile split and under a spatial block split. Sections~\ref{sec:classical}--\ref{sec:yolo} document the four methods; Sect.~\ref{sec:comparison} presents the shared-protocol comparison and the split-protocol study; Sect.~\ref{sec:southpole} deploys the pipelines on the unlabelled lunar south pole as a transfer diagnostic. The central methodological finding is that on planetary datasets with extreme prevalence structure, the evaluation protocol --- split geometry, metric choice, threshold calibration --- moves reported numbers by more than any architectural decision we tested, and must be fixed before architectures are compared.
